\section{Virtualisierung}
\textbf{Virtualisierung} ist ein schon lange bestehendes Konzept. Dabei wird auf einem
physischen Computer per Software eine andere, \emph{isolierte Umgebung} abstrahiert.
Bekannt ist das Konzept beispielsweise durch \textit{Oracle VirtualBox}. Als Abstraktion
von Hardware-Ressourcen durch Software ist Virtualisierung zugleich der
\emph{Grundstein} für Cloud-Technologie: Erst sie ermöglicht \textit{Resource Pooling}, denn
sonst könnten keine voneinander getrennten Umgebungen auf einem Server
existieren.
\subsection{Wie nutzt eine Cloud Virtualisierung?}
\subsubsection{Virtuelle Maschinen (Virtual Machine)}
Bei einer \textbf{VM} (kurz für Virtuelle Maschine) wird eine \emph{eigenständige, isolierte}
Instanz auf einem physischen System abstrahiert. Dadurch kann man
beispielsweise auf einem Windows-PC ein Linux-Betriebssystem innerhalb einer VM
ausführen. In der Cloud wird dies für \textit{Resource Pooling} genutzt, also das
Aufteilen eines physischen Servers in mehrere kleinere, unabhängige Instanzen.
So werden \textit{Infrastructure as a Service (IaaS)} und dynamische Skalierung erst
ermöglicht.
\subsubsection{Virtuelles Netzwerk}
Ein \textbf{virtuelles Netzwerk} ermöglicht die Kommunikation zwischen VMs sowie mit
externen Systemen. Durch \textit{Peering} – das direkte Verbinden zweier virtueller
Netzwerke – können Services untereinander kommunizieren, ohne den Umweg über
das öffentliche Internet nehmen zu müssen. Über \textit{VPNs} (Virtual Private Network)
lässt sich zudem eine verschlüsselte Direktverbindung zu
\textit{On-Premises}-Rechenzentren herstellen, um sicher mit Cloud-Services zu
kommunizieren.
\subsubsection{Speichervirtualisierung}
Bei der \textbf{Speichervirtualisierung} werden mehrere physische Speichermedien per
Software zu einem einzigen logischen Speicherpool zusammengefasst und
abstrahiert. Das ermöglicht beispielsweise Dienste wie \textit{Google Drive} und \textit{iCloud}
oder – im Enterprise-Kontext – \textit{Azure Managed Disks}.

\subsubsection{Hypervisor}
Ein \textbf{Hypervisor} ist die Software, die Arbeitsspeicher, Rechenleistung und
Speicher in Pools zusammenfasst. Dadurch lassen sich auf einer physischen
Maschine viele VMs erstellen. Der Hypervisor stellt den einzelnen virtuellen
Maschinen die zugewiesenen Ressourcen bereit und verwaltet deren Planung im
Verhältnis zu den physischen Ressourcen. Die Ausführung übernimmt nach wie vor
\emph{die physische Hardware}: Die CPU führt weiterhin die CPU-Befehle aus, wie sie
etwa von den VMs angefordert werden, während der Hypervisor die Planung
übernimmt.

\begin{table}[h]
    \centering
    \begin{tabular}{|l|p{5cm}|p{5cm}|}
        \hline
        \textbf{Merkmal} & \textbf{Typ 1 (Bare-Metal)} & \textbf{Typ 2 (Gehostet)} \\
        \hline
        Ausführung & Direkt auf der Hardware & Auf einem Host-Betriebssystem \\
        \hline
        Host-OS & Keins, ersetzt das OS & Benötigt ein Host-OS \\
        \hline
        Ressourcenzuweisung & Direkt durch den Hypervisor & Über das Host-OS \\
        \hline
        Beispiele & \textit{VMware ESXi}, \textit{KVM} & \textit{Oracle VirtualBox}, \textit{VMware Workstation} \\
        \hline
    \end{tabular}
    \caption{Vergleich Hypervisor Typ~1 und Typ~2}
    \label{tab:hypervisor}
\end{table}


